% =============================================================================
% Methodology section for the H2 paper.
% Written as a drop-in LaTeX section using only standard packages:
%   amsmath, amssymb, amsthm, algorithm, algpseudocode
% Put it inside any TMLR template as \section{Methodology}.
%
% CORRECTED VERSION. Differences from the previous draft:
%   - Proposition 1 now states the exact (d-1)/2 parameterization, and a
%     following Remark documents the Beta(d/2, d/2) asymptotic convention used
%     in the code, consistent with the TurboQuant lineage.
%   - Equation for the dequantization now uses M, not M^T, matching the code.
%   - Added explicit subsections describing the baseline (no-rotation + codebook)
%     and RTN-absmax (canonical LLM) baselines used in our experiments.
%   - Lloyd-Max construction now documents that iteration proceeds in [0, 1]
%     and is shifted to [-1, 1] at the end, matching the implementation.
%   - Codebook storage arithmetic corrected.
% =============================================================================

\section{Methodology}
\label{sec:method}

We describe H2, a post-training weight quantization pipeline that combines
random orthogonal rotation with a scalar codebook matched to the coordinate
distribution that rotation produces. The method is training-free, requires no
calibration data, and runs in under a minute on a single CPU core for
ViT-B/16. This section derives H2 from first principles, states each
assumption precisely, enumerates the baselines against which we compare it,
and identifies the architectural regime in which the approach applies.

\subsection{Problem setup and notation}
\label{sec:method-setup}

Consider a pretrained neural network with $L$ weight layers indexed by
$\ell \in \{1, \ldots, L\}$. For each layer we flatten the weight tensor to a
matrix $W_\ell \in \mathbb{R}^{N_\ell \times d_\ell}$ using the following
conventions. A two-dimensional convolution with kernel shape
$(C_\text{out}, C_\text{in}, k_H, k_W)$ flattens to a matrix of shape
$(C_\text{out},\, C_\text{in}\, k_H\, k_W)$, where each row corresponds to
one output filter. A fully-connected layer with shape
$(\text{out}, \text{in})$ is already in two-dimensional form. A standard
multi-head attention module stores the query, key, and value projections as
a single parameter \texttt{in\_proj\_weight} of shape $(3d, d)$, and we
treat this parameter as a linear weight with fan-in $d$. Throughout,
$N_\ell$ denotes the number of output rows of layer $\ell$ and $d_\ell$
denotes its fan-in.

Given a bit budget $b$, the quantization problem is to produce an
approximation $\widehat{W}_\ell$ that uses at most $b$ bits per element,
plus a small amount of per-row floating-point metadata, and that preserves
the task-level accuracy of the network. We write row $i$ of layer $\ell$ as
$w_{\ell,i} \in \mathbb{R}^{d_\ell}$.

\subsection{Motivation for a matched codebook}
\label{sec:method-motivation}

A uniform quantizer with $2^b$ levels equally spaced on the interval
$[-1, 1]$ has step size $2/2^b$. Applied directly to a pretrained weight
matrix it suffers two well-known problems. The empirical weight
distribution concentrates sharply near zero with rare large-magnitude
outliers, so uniformly spaced levels waste resolution on regions of
near-zero probability mass. At the same time, per-channel absmax scaling of
the form $q = \text{round}(x \cdot L / s)$ with $s = \max |x|$ and level
count $L = 2^{b-1} - 1$ is dominated by outliers, which compresses the
effective dynamic range of typical weights.

The QuaRot approach of \citet{ashkboos2024quarot} addresses the outlier
problem by applying a Hadamard rotation before per-channel absmax
quantization. The rotation redistributes outlier information uniformly
across coordinates, which shrinks the per-row dynamic range. The
quantization step, however, still uses a uniform grid, which does not match
the post-rotation coordinate distribution. We characterize that
distribution precisely, and then construct a codebook optimal for it.

\subsection{Coordinate distribution on the unit sphere}
\label{sec:method-distribution}

The method rests on one classical fact from spherical geometry.

\begin{proposition}[Sphere coordinate distribution]
\label{prop:sphere}
Let $x$ be drawn uniformly from the unit sphere
$S^{d-1} \subset \mathbb{R}^d$ with $d \geq 2$. Then each coordinate $x_j$
for $j \in \{1, \ldots, d\}$ has marginal density on the interval
$[-1, 1]$ given by
\begin{equation}
f_d(u) \;=\;
   \frac{1}{B\!\left(\tfrac{d-1}{2}, \tfrac{d-1}{2}\right) \cdot 2^{d-2}}
   \cdot (1 - u^2)^{(d-3)/2},
\label{eq:spherical-density}
\end{equation}
which is the density of the rescaled Beta random variable
\begin{equation}
x_j \;\sim\; 2\,\mathrm{Beta}\!\left(\tfrac{d-1}{2}, \tfrac{d-1}{2}\right) - 1.
\label{eq:spherical-rescaled}
\end{equation}
\end{proposition}

\begin{proof}
Let $g \sim \mathcal{N}(0, I_d)$. The rotational invariance of the Gaussian
density implies that $x = g / \|g\|_2$ is uniform on $S^{d-1}$. Consider
the squared first coordinate
\begin{equation}
x_1^2 \;=\; \frac{g_1^2}{\sum_{i=1}^d g_i^2}
          \;=\; \frac{g_1^2}{g_1^2 + \sum_{i=2}^d g_i^2}.
\label{eq:sq-ratio}
\end{equation}
Since $g_i^2 \sim \chi^2_1$ are independent and
$\sum_{i=2}^d g_i^2 \sim \chi^2_{d-1}$, the ratio is distributed as
$x_1^2 \sim \mathrm{Beta}(\tfrac{1}{2}, \tfrac{d-1}{2})$ by the standard
Beta-from-ratio-of-chi-squares identity. Applying the change of variables
$u = \pm\sqrt{v}$ for $v \sim \mathrm{Beta}(\tfrac{1}{2}, \tfrac{d-1}{2})$
and symmetrizing over the sign yields the density
$f_d(u) \propto (1-u^2)^{(d-3)/2}$ on $[-1, 1]$. The normalization follows
by matching this density to the rescaled Beta distribution
$\mathrm{Beta}(\tfrac{d-1}{2}, \tfrac{d-1}{2})$ on $[-1, 1]$, whose density
is obtained from the Beta density on $[0, 1]$ by the linear change of
variable $y = (u+1)/2$ with Jacobian $1/2$.
\end{proof}

\paragraph{Explanation of Equation \ref{eq:spherical-density}.}
The factor $(1 - u^2)^{(d-3)/2}$ controls the shape of the density and tends
to zero as $|u| \to 1$, pushing probability toward the origin as $d$ grows.
The constants in front, namely $1 / B((d-1)/2, (d-1)/2)$ and $1 / 2^{d-2}$,
together ensure that the density integrates to one on $[-1, 1]$. The Beta
function
$B(\alpha, \beta) = \int_0^1 y^{\alpha-1}(1-y)^{\beta-1}\, dy$ normalizes the
Beta density on $[0, 1]$, and the factor $1 / 2^{d-2}$ accounts for the
rescaling from $[0, 1]$ to $[-1, 1]$.

\begin{remark}[Asymptotic convention]
\label{rem:asymptotic}
Our implementation follows the TurboQuant paper \citep{turboquant} and its
reference code, which use the convention
$x_j \sim 2\,\mathrm{Beta}(d/2, d/2) - 1$ throughout. This is a large-$d$
approximation of the exact result in Equation~\ref{eq:spherical-rescaled},
and we adopt it here for direct correspondence with that code path. Three
properties make the approximation harmless in our regime. First, the
variance of $\mathrm{Beta}((d-1)/2, (d-1)/2)$ rescaled to $[-1, 1]$ equals
$1/d$, while that of $\mathrm{Beta}(d/2, d/2)$ equals $1/(d+1)$, and the
two differ by $1/(d(d+1)) = O(1/d^2)$. Second, the KL divergence between
these two distributions is $O(1/d)$ by a direct digamma expansion, so by
Pinsker's inequality the total variation distance is $O(1/\sqrt{d})$.
Third, the Lloyd--Max codebook constructed under either density has
distortion that is within a factor $1 + O(1/d)$ of the optimum under the
other, because Lloyd--Max distortion is a continuous functional of the
source density in the Hellinger metric. Our implementation therefore uses
$\mathrm{Beta}(d/2, d/2)$ for both the Lloyd--Max codebook of
Section~\ref{sec:method-codebook} and the Kolmogorov--Smirnov validation
test of Section~\ref{sec:experiments}, and we report all theoretical
statements using the exact Beta$((d-1)/2, (d-1)/2)$ parameterization. The
accuracy loss from this mismatch is negligible when $d \gtrsim 100$, which
covers every layer to which we apply H2 in the experiments.
\end{remark}

\paragraph{Consequence at large $d$.}
Under either parameterization the density has mean zero by symmetry. The
variance equals $1/(d+1)$ under the code's Beta$(d/2, d/2)$ convention and
equals $1/d$ under the exact Beta$((d-1)/2, (d-1)/2)$ convention. Both
expressions agree at leading order and give standard deviation of order
$1/\sqrt{d}$. For $d = 768$, which is the fan-in of a ViT-B/16 attention
projection, the standard deviation is approximately $0.036$ under either
convention, and more than $95\%$ of the probability mass lies inside
the interval $[-0.07, +0.07]$. A codebook that treats values near $\pm 1$ as
equally important as values near zero wastes most of its levels on
probability mass that does not exist.

\subsection{Isotropization by random rotation}
\label{sec:method-rotation}

Pretrained weight rows do not themselves sit uniformly on the sphere.
Training imparts directional structure that encodes the network's learned
features. We approximately remove that structure by a two-step
transformation applied to each layer independently.

\paragraph{Step 1. Row normalization.}
For each row $w_i$ of $W$ we compute the Euclidean norm $n_i = \|w_i\|_2$
and the normalized row $\widetilde{w}_i = w_i / n_i$. The normalized row
lies on $S^{d-1}$, but it reflects whatever directional structure remains
after training. The norm $n_i$ is stored as floating-point metadata, at a
cost of $4N$ bytes per layer. Our implementation factors this step with
Step~2 below for numerical convenience, applying the rotation to the
unnormalized row first and then normalizing the rotated output. The two
orderings are mathematically equivalent because rotation is an isometry,
so $\|M w_i\| = \|w_i\| = n_i$ regardless of whether normalization
precedes or follows the rotation.

\paragraph{Step 2. Random orthogonal rotation.}
We sample a random orthogonal matrix $M \in O(d)$ uniformly from the Haar
measure and form $u_i = M\, \widetilde{w}_i$, viewing $\widetilde{w}_i$ as a
column vector. In matrix form the operation on all rows simultaneously is
\begin{equation}
U_{\text{matrix}}
  \;=\; \widetilde{W}_{\text{matrix}}\, M^\top,
\label{eq:forward-matrix}
\end{equation}
where $\widetilde{W}_{\text{matrix}}$ has row $i$ equal to
$\widetilde{w}_i^\top$, and the factor $M^\top$ appears because post-
multiplying by $M^\top$ is equivalent to pre-multiplying each column-vector
row by $M$. Because the uniform measure on $S^{d-1}$ is rotation-invariant,
a uniformly random rotation disperses any directional structure the
original row possessed. In the idealized case where $\widetilde{w}_i$ is
uniform on the sphere, $u_i$ remains uniform and Proposition~\ref{prop:sphere}
applies exactly. For pretrained weights the empirical coordinate
distribution of $u_i$ approaches $f_d$ as $d$ grows, a convergence we
verify empirically using Kolmogorov--Smirnov statistics in
Section~\ref{sec:experiments}.

We sample $M$ via one of two constructions depending on $d$. When $d$ is a
power of two we use the Structured Randomized Hadamard Transform of the
form $M = D_1\, H\, D_2$, where $H$ is the $d \times d$ Hadamard matrix
and $D_1, D_2$ are independent diagonal sign matrices with entries drawn
uniformly from $\{-1, +1\}$. This construction has $O(d)$ memory and
$O(d \log d)$ application cost. When $d$ is not a power of two we
materialize a dense orthogonal matrix by taking the $Q$ factor from the
QR decomposition of a $d \times d$ matrix with i.i.d. standard Gaussian
entries and applying the Mezzadri \citep{mezzadri2007} sign correction
$Q \leftarrow Q\, \mathrm{diag}(\mathrm{sign}(\mathrm{diag}(R)))$ to
recover uniform Haar measure on $O(d)$, at $O(d^2)$ memory and
application cost. Only the seed used to
sample $M$ is persisted. The matrix itself is regenerated at quantization
time and again at dequantization, then discarded.

\subsection{Lloyd--Max codebook for the Beta density}
\label{sec:method-codebook}

Given a fan-in $d$ and a bit budget $b$, we construct a scalar codebook of
$K = 2^b$ levels that minimizes expected squared reconstruction error under
the density $f_d$ from Remark~\ref{rem:asymptotic}, which we identify with
$\mathrm{Beta}(d/2, d/2)$ rescaled to $[-1, 1]$. A scalar codebook consists
of $K$ decision boundaries
$-1 = t_0 < t_1 < \cdots < t_K = 1$ and $K$ centroids $c_1, \ldots, c_K$
with $c_k \in [t_{k-1}, t_k]$. The expected squared reconstruction error
is
\begin{equation}
\mathcal{D}(\{t_k\}, \{c_k\})
  \;=\; \sum_{k=1}^{K} \int_{t_{k-1}}^{t_k}
          (u - c_k)^2\, f_d(u)\, du.
\label{eq:distortion}
\end{equation}

The classical Lloyd--Max necessary conditions
\citep{lloyd1982, max1960} for a stationary point of
Equation~\ref{eq:distortion} are
\begin{equation}
c_k \;=\; \mathbb{E}[U \mid t_{k-1} < U < t_k]
       \;=\; \frac{\int_{t_{k-1}}^{t_k} u\, f_d(u)\, du}
                  {\int_{t_{k-1}}^{t_k} f_d(u)\, du},
\label{eq:centroid-update}
\end{equation}
\begin{equation}
t_k \;=\; \tfrac{1}{2}\left( c_k + c_{k+1} \right),
       \quad k \in \{1, \ldots, K-1\}.
\label{eq:boundary-update}
\end{equation}

\paragraph{Explanation of Equations \ref{eq:centroid-update}
and \ref{eq:boundary-update}.}
Equation~\ref{eq:centroid-update} states that the optimal centroid within a
given bin equals the conditional mean of the distribution on that bin. This
is the unique minimizer of the squared-error loss restricted to the bin.
Equation~\ref{eq:boundary-update} states that the optimal decision
boundary between two adjacent centroids lies at their midpoint. A point
exactly halfway between two centroids incurs the same reconstruction error
under either, so this is the tie-breaking location.

\paragraph{Numerical construction in $[0, 1]$.}
Our implementation performs the Lloyd--Max iteration on the equivalent
Beta density over $[0, 1]$ and shifts to $[-1, 1]$ at the end. Specifically
we initialize the boundaries to the Beta quantiles on $[0, 1]$,
$t_k^{(0)} = F_d^{-1}(k/K)$ where $F_d$ is the CDF of $\mathrm{Beta}(d/2,
d/2)$ on $[0, 1]$, with $t_0^{(0)} = 0$ and $t_K^{(0)} = 1$. We then
alternate the updates in Equations~\ref{eq:centroid-update}
and~\ref{eq:boundary-update} until the elementwise update of the
boundary vector falls below $10^{-10}$ in absolute value, where
numerical integration uses the trapezoidal rule with one thousand sample
points per bin. Convergence typically requires fewer than one hundred
iterations. At
termination we apply the affine shift $t_k \leftarrow 2 t_k - 1$ and
$c_k \leftarrow 2 c_k - 1$ so that the codebook lives on the target
interval $[-1, 1]$. Working in $[0, 1]$ uses the native support of the
Beta density and avoids a change-of-variables inside the integrands, while
yielding an equivalent codebook because the transformation is affine.

\paragraph{Convergence.}
The density $f_d$ is log-concave for $d \geq 3$, and classical results
\citep{fleischer1964, kieffer1983} guarantee that Lloyd--Max converges to
the unique global optimum of Equation~\ref{eq:distortion} under
log-concavity. For $d = 2$ the density $f_2$ is the arcsine density, which
is not log-concave. This corner case does not arise for any layer in the
architectures we study, so the log-concavity condition is satisfied
throughout.

\paragraph{Codebook caching.}
The codebook is determined entirely by the pair $(d, b)$ and does not depend
on the specific weights. We precompute one codebook per distinct fan-in and
share it across all layers that have matching $d$ and $b$. A typical vision
or language model contains fewer than ten distinct fan-ins across all of
its layers, so the total codebook overhead is negligible compared to the
weight storage.

\subsection{The H2 quantization algorithm}
\label{sec:method-pipeline}

We now combine the pieces of Sections \ref{sec:method-rotation}
and~\ref{sec:method-codebook} into the full algorithm. Given a weight
matrix $W \in \mathbb{R}^{N \times d}$, a bit budget $b$, and a rotation
seed $s$, the procedure produces a quantized reconstruction $\widehat{W}$
along with per-row norms and the rotation seed.

\begin{algorithm}[t]
\caption{H2 post-training weight quantization, applied per layer.}
\label{alg:h2}
\begin{algorithmic}[1]
\Require weight matrix $W \in \mathbb{R}^{N \times d}$, bit budget $b$,
         rotation seed $s$
\Ensure  quantized codes $\widehat{I} \in \{0, \ldots, 2^b-1\}^{N \times d}$,
         norms $n \in \mathbb{R}^N$, seed $s$
\State $n_i \gets \|w_i\|_2$ for each row $i$
\Comment{L2 row norms}
\State $\widetilde{W} \gets \mathrm{diag}(n)^{-1}\, W$
\Comment{row-normalize}
\State $M \gets \textsc{SampleOrthogonal}(d, s)$
\Comment{SRHT or dense}
\State $U \gets \widetilde{W}\, M^\top$
\Comment{rotate row-wise}
\State $(\{t_k\}, \{c_k\}) \gets \textsc{LloydMax}(f_d, K{=}2^b)$
\Comment{cached}
\For{$i \in \{1, \ldots, N\},\ j \in \{1, \ldots, d\}$}
  \State $\widehat{I}_{i,j} \gets k$ such that $t_{k-1} \leq U_{i,j} < t_k$
\EndFor
\State \Return $\widehat{I}, n, s$
\end{algorithmic}
\end{algorithm}

Dequantization at inference time proceeds by recovering $M$ from the seed,
reconstructing the quantized matrix $\widehat{U}$ from the codes
$\widehat{I}$ via $\widehat{U}_{i,j} = c_{\widehat{I}_{i,j}}$, and then
computing
\begin{equation}
\widehat{W}
  \;=\; \mathrm{diag}(n)\, \widehat{U}\, M.
\label{eq:dequant}
\end{equation}

\paragraph{Explanation of Equation \ref{eq:dequant}.}
The two operations on the right-hand side invert the two steps applied in
Algorithm~\ref{alg:h2}. The forward rotation in matrix form was
$U_{\text{matrix}} = \widetilde{W}_{\text{matrix}}\, M^\top$ as in
Equation~\ref{eq:forward-matrix}. Its inverse is obtained by right-
multiplying by $(M^\top)^{-1} = M$, since $M$ is orthogonal and therefore
$M^\top M = I$. This gives
$\widetilde{W}_{\text{matrix}} = U_{\text{matrix}}\, M$, which when combined
with the row scaling $\mathrm{diag}(n)$ produces the full inverse shown
above. The resulting matrix $\widehat{W}$ replaces $W$ in the network, and
the forward pass proceeds using $\widehat{W}$ exactly as the unquantized
network uses $W$.

\subsection{Baselines in the training-free regime}
\label{sec:method-baselines}

Our experiments compare H2 against three baselines, each of which is also
training-free and requires no calibration data. We describe them formally
here so that every method appearing in our result tables has a
corresponding algorithmic specification.

\paragraph{RTN-absmax.}
Per-channel symmetric round-to-nearest quantization with no rotation. This
is the canonical baseline used in the large language model quantization
literature, including GPTQ \citep{frantar2022gptq}, AWQ
\citep{lin2023awq}, and QuaRot \citep{ashkboos2024quarot}. For each row
$w_i$ of $W$ we set $s_i = \max_j |w_{i,j}|$, choose an integer range
$L = 2^{b-1} - 1$, and quantize each element to
$\widehat{I}_{i,j} = \mathrm{round}(w_{i,j}\, L / s_i)$ clamped to
$[-L, L]$. Dequantization is
$\widehat{w}_{i,j} = \widehat{I}_{i,j}\, s_i / L$. The only per-row metadata
is the scale $s_i$. RTN-absmax has no rotation step and no distributional
assumption on the weights.

\paragraph{QuaRot-RTN.}
Rotation followed by RTN-absmax. We sample the same kind of random
orthogonal matrix $M$ as in H2 and apply the forward rotation
$U = \widetilde{W}\, M^\top$. For H2 the rotation is applied after row
L2 normalization, whereas for QuaRot-RTN we follow the original
prescription and apply the rotation directly to the raw row, then set the
scale as $s_i = \max_j |U_{i,j}|$. Quantization and dequantization proceed
as in RTN-absmax, operating on the rotated matrix $U$. The recovered
matrix is $\widehat{W} = \widehat{U}\, M$. We compare against this
training-free variant of QuaRot for an apples-to-apples isolation of the
codebook choice. The full QuaRot paper additionally uses GPTQ-style
calibration-based weight refinement, which is orthogonal to the rotation
and codebook design and can be applied on top of either QuaRot-RTN or H2.

\paragraph{Baseline with codebook but no rotation.}
To isolate the contribution of the rotation step, we also run a variant
that applies row L2 normalization and the Lloyd--Max Beta codebook but
omits the rotation. This is Algorithm~\ref{alg:h2} with step 3 replaced by
$M = I$. The resulting quantizer is straightforward to derive
theoretically, because unrotated normalized weights do not have the
Beta density, and so the codebook is mismatched to the actual distribution.
We report the performance of this variant in our result tables to make
visible the independent contribution of rotation, separately from the
contribution of a matched codebook.

\subsection{Storage accounting}
\label{sec:method-storage}

Each quantized layer stores three components. The quantized codes occupy
$N\, d\, b / 8$ bytes. The per-row norms for H2 and per-row scales for
RTN-absmax and QuaRot-RTN require $4N$ bytes in floating-point precision.
The rotation seed is a single four-byte integer. The codebook itself
consists of $K + 1 = 2^b + 1$ boundaries and $K = 2^b$ centroids, for a
total of $4\,(2 \cdot 2^b + 1)$ bytes, and it is shared across all layers
with the same fan-in and bit budget.

For a concrete example, ViT-B/16 at $b=4$ contains approximately $65$
million quantizable weights across $38$ linear layers plus multi-head
attention projections. The codes occupy about $32.5$ MB, the row norms
occupy under $200$ KB, non-quantizable parameters such as layer
normalization gains and positional embeddings occupy approximately $6$ MB
in floating-point precision, and the codebook occupies one hundred and
thirty-two bytes. The total effective size is therefore approximately
$39$ MB, compared to the $330$ MB floating-point baseline, which
corresponds to a compression ratio near $8.5$. The exact ratio reported
in our result tables is $7.8$, with the discrepancy accounted for by the
fact that our reporting includes all non-quantizable parameters in the
FP32 reference and conservatively keeps them at FP32 in the compressed
budget, whereas this back-of-envelope figure assumes FP16 for all
non-quantizable parameters.

\subsection{Handling of multi-head attention}
\label{sec:method-mha}

Standard \texttt{nn.MultiheadAttention} modules in PyTorch expose the
combined query, key, and value projection as a single parameter named
\texttt{in\_proj\_weight} of shape $(3d, d)$. This parameter sits outside
any \texttt{nn.Linear} submodule, and a naive iteration that visits only
convolutional and linear modules does not encounter it. For ViT-B/16 this
omission affects $21.5$ million parameters, which is $33\%$ of the
quantizable weight budget. Leaving those parameters at floating-point
precision while quantizing the rest of the model produces misleadingly
strong compressed accuracy at the cost of substantially understating the
true compression ratio.

We augment the per-layer iteration with a dedicated pass over
\texttt{nn.MultiheadAttention} modules that runs after the main iteration
completes. For each such module we apply Algorithm~\ref{alg:h2} to
\texttt{in\_proj\_weight} as if it were a linear weight with fan-in $d$
equal to the embedding dimension. The \texttt{out\_proj} submodule is
already an \texttt{nn.Linear} and is handled by the standard pass. Causal
language models built on the Llama or GPT architectures expose the four
attention projections as separate \texttt{nn.Linear} modules, so the
multi-head attention pass is specific to vision-transformer
implementations that use PyTorch's default module.

\subsection{Why H2 beats uniform after rotation}
\label{sec:method-whybeats}

QuaRot-RTN and H2 apply the same rotation and both use per-row scale
information. The methods differ only in the codebook and in the choice of
scale. Comparing them head-to-head therefore isolates the effect of
matching the codebook to the post-rotation distribution.

Consider $b = 2$, which leaves four levels per coordinate. Under
QuaRot-RTN the four levels are placed at $\{-1, -1/3, +1/3, +1\}$
multiplied by the per-row absmax, which for typical weight rows after
rotation equals three to four standard deviations of the coordinate
distribution. The three interior transitions at $\pm 1/3$ and zero are
therefore far from where the density concentrates, because the standard
deviation under $f_d$ is of order $1/\sqrt{d}$ relative to the unit-norm
scale. Under H2, Lloyd--Max places the four centroids at the conditional
means of the four quartiles of $f_d$. At large $d$ the density is
approximately Gaussian with standard deviation $\sigma \approx
1/\sqrt{d}$, and the classical Lloyd--Max solution for a unit-variance
Gaussian source at four levels \citep{max1960} places the centroids at
$\pm 0.4528\sigma$ and $\pm 1.510\sigma$. Translated to our scale this
gives centroids at approximately $\pm 0.4528/\sqrt{d}$ and
$\pm 1.510/\sqrt{d}$, which concentrates the representational budget in
the high-density region near zero. The gap between the two arrangements
shrinks as $b$ grows, because with enough levels both codebooks cover the
distribution adequately.

Classical analyses of high-resolution scalar quantization
\citep{bennett1948, panter1951, gershogray1992} give the asymptotic
distortion of the optimal (Lloyd--Max) quantizer for a source with
density $f$ at $K$ levels as
\begin{equation}
\mathcal{D}^\star_K \;\approx\; \frac{1}{12\, K^2}
    \left( \int f(x)^{1/3}\, dx \right)^3,
\label{eq:bennett}
\end{equation}
which is Bennett's integral. The constant in front depends on the shape
of the density through the functional $\big(\int f^{1/3}\big)^3$, not on
the second moment. Two densities with the same variance can have very
different Bennett integrals if their shapes differ.

For our rotated coordinates we can substitute $u = v/\sqrt{d}$ and apply
the Gaussian approximation to $f_d$, which holds to leading order for
large $d$ because $\mathrm{Beta}((d-1)/2, (d-1)/2)$ rescaled to $[-1, 1]$
concentrates on scale $1/\sqrt{d}$ around zero. This gives
\begin{equation}
\int f_d(u)^{1/3}\, du
\;\approx\;
\left(\frac{d}{2\pi}\right)^{1/6} \int \exp\!\left(-\frac{d u^2}{6}\right) du
\;=\;
C \cdot d^{-1/3},
\label{eq:bennett-beta}
\end{equation}
so Bennett's formula yields $\mathcal{D}^\star_K = O(K^{-2}/d)$ for H2.
The comparison to QuaRot-RTN requires care because QuaRot uses per-row
adaptive scaling rather than a fixed grid. Each row is quantized
uniformly on $[-s_i, +s_i]$ where $s_i = \max_j |U_{i,j}|$ is the per-row
absolute maximum. Under the Gaussian approximation to $f_d$, the $d$
rotated coordinates of a unit-norm row are approximately
$\mathcal{N}(0, 1/d)$, and the standard extreme-value bound gives
$s_i \sim \sqrt{2 \log d / d}$ almost surely for large $d$. The step
size of QuaRot's uniform quantizer is therefore
$\Delta = 2 s_i / K$, and the high-resolution distortion at mid-bin is
\begin{equation}
\mathcal{D}^{\text{QuaRot}}_K
   \;\approx\; \frac{\Delta^2}{12}
   \;=\; \frac{s_i^2}{3 K^2}
   \;=\; O\!\left( \frac{\log d}{d\, K^2} \right).
\label{eq:quarot-distortion}
\end{equation}

\paragraph{Explanation of Equation \ref{eq:quarot-distortion}.}
The adaptive absmax scaling closes most of the gap against a fixed-grid
uniform quantizer, and what remains is a $\log d$ mismatch between the
absmax scale $s_i \sim \sqrt{\log d / d}$ and the natural density scale
$1/\sqrt{d}$ of the bulk. QuaRot's 16 or 256 levels at $b \in \{4, 8\}$
fall across the absmax range, but most of the probability mass lives
within a fraction $1/\sqrt{\log d}$ of that range, and the unused outer
levels are a quantified waste.

The ratio of the two distortion rates is
\begin{equation}
\frac{\mathcal{D}^{\text{QuaRot}}_K}{\mathcal{D}^\star_K}
   \;=\; O(\log d),
\label{eq:ratio-logd}
\end{equation}
which is the leading-order advantage of a density-matched codebook over
adaptive-uniform quantization on the same rotated coordinates. This is a
much weaker claim than the $O(d)$ advantage of density-matched over
fixed-grid uniform, because QuaRot's absmax scaling already does most of
the work.

\paragraph{Caveat at low bit width.}
Equations~\ref{eq:bennett}, \ref{eq:bennett-beta}
and~\ref{eq:quarot-distortion} are high-resolution asymptotics that
become numerically tight as $K \to \infty$. At $K = 4$, which is the
2-bit extreme compression regime in our experiments, these
approximations are not numerically accurate. What the asymptotic
analysis does correctly predict is the direction and fan-in dependence
of the gap between density-matched and adaptive-uniform quantization, a
prediction we confirm empirically in Section~\ref{sec:experiments}.
The empirical head-to-head gap at $K = 4$ is frequently much larger
than $\log d$ would suggest, because at low resolution the placement of
a handful of levels relative to the bulk of the density dominates the
distortion, and QuaRot's outer levels at $\pm s_i$ are far from where
probability lives.

\subsection{Architectural boundary}
\label{sec:method-boundary}

Proposition~\ref{prop:sphere} applies exactly when $\widetilde{w}$ is
uniform on $S^{d-1}$ and applies approximately when a random rotation
applied to a structured $\widetilde{w}$ produces coordinates close to
uniform. Both parts of this argument become weak when $d$ is small, for
two distinct reasons.

The target density $f_d$ itself loses concentration at small $d$. At
$d = 9$, which is the fan-in of a three-by-three depthwise convolution,
the code's $\mathrm{Beta}(4.5, 4.5)$ convention gives a standard deviation
of approximately $0.316$ on $[-1, 1]$. This is more peaked than the
uniform density on the same interval, whose standard deviation is
approximately $0.577$, but the concentration is not dramatic. The
Lloyd--Max levels for $f_9$ are only modestly closer to zero than
uniformly spaced levels, so the codebook advantage over uniform is small.

A random rotation in nine dimensions cannot produce the asymptotic
isotropization that holds when $d$ is large. The orthogonal group
$O(9)$ is rich enough to mix coordinates, but with only nine values per
filter the empirical coordinate distribution is sensitive to the specific
sample and retains visible structure from the original weights. In
addition, L2 normalization over only nine values is statistically noisy,
which further degrades the match to the target density.

Taken together, these two factors explain the one architecture in our
study where H2 fails, namely MobileNet-V2. Depthwise convolutions with
$d = 9$ do not benefit from the Beta codebook, and the per-row L2
normalization adds noise without improving the distributional match. The
depthwise-separable blocks in EfficientNet-B0 and ConvNeXt-Tiny are less
severe cases, because those networks interleave large-fan-in pointwise
convolutions with small-fan-in depthwise layers. The pointwise layers
carry enough of the quantization budget that H2 still provides an overall
benefit, even though the depthwise layers are in the same flat-density
regime as MobileNet-V2.

The practical rule from this analysis is that H2 should be preferred when
$d \gtrsim 100$, and that a uniform codebook with rotation, which is to
say QuaRot-RTN, should be preferred when $d \lesssim 20$. The transition
region between these two regimes contains most of the depthwise-separable
architectures in common use, and choosing the right codebook per layer is
an obvious avenue for future work.

\subsection{Computational cost}
\label{sec:method-cost}

H2 introduces no overhead at inference time beyond the one-time
dequantization step when the model is loaded. The forward pass operates on
the reconstructed $\widehat{W}$ exactly as the original model operates on
$W$, and on hardware with low-precision arithmetic support the natural
speedup from smaller weights applies without further modification.

The quantization step itself is dominated by the construction and
application of the rotation matrix. For a layer with fan-in $d$ the SRHT
construction costs $O(d \log d)$ per row, and the dense orthogonal
construction costs $O(d^2)$ per row. The Lloyd--Max iteration is amortized
across all layers with matching fan-in, so its cost is negligible. For
ViT-B/16 the complete quantization at four bits, including rotation,
codebook construction, quantization, and dequantization for all layers,
takes under thirty seconds on a single CPU core. No GPU is required.

\subsection{Summary}
\label{sec:method-summary}

H2 is the composition of three well-understood ingredients. We normalize
each row of the weight matrix to the unit sphere, apply a uniformly
random orthogonal rotation to the fan-in axis, and quantize the rotated
coordinates against a Lloyd--Max codebook tuned to the Beta density that
arises on $S^{d-1}$. The first two steps produce coordinates whose
distribution is known analytically in the large-$d$ limit, and the third
step is the minimum-distortion codebook for that distribution. All three
steps are training-free and deterministic given the rotation seed. The
method's effective regime, which is layers with $d \gtrsim 100$, is
characterized by the same geometric argument that motivates the
construction, and the architectural boundary where the method fails is
explained by the same argument in reverse.
